
\documentclass{article}
\usepackage{colortbl}
\usepackage{makecell}
\usepackage{multirow}
\usepackage{supertabular}

\begin{document}

\newcounter{utterance}

\centering \large Interaction Transcript for game `hot\_air\_balloon', experiment `air\_balloon\_survival\_hr\_complexity\_hard', episode 2 with glm{-}5.2.
\vspace{24pt}

{ \footnotesize  \setcounter{utterance}{1}
\setlength{\tabcolsep}{0pt}
\begin{supertabular}{c@{$\;$}|p{.15\linewidth}@{}p{.15\linewidth}p{.15\linewidth}p{.15\linewidth}p{.15\linewidth}p{.15\linewidth}}
    \# & $\;$A & \multicolumn{4}{c}{Game Master} & $\;\:$B\\
    \hline

    \theutterance \stepcounter{utterance}  
    & & \multicolumn{4}{p{0.6\linewidth}}{
        \cellcolor[rgb]{0.9,0.9,0.9}{
            \makecell[{{p{\linewidth}}}]{
                \texttt{\tiny{[P1$\langle$GM]}}
                \texttt{Sudjelujete u suradničkoj pregovaračkoj igri.} \\
\\ 
\texttt{Zajedno s drugim sudionikom morate se dogovoriti o jednom skupu predmeta koji će se zadržati. Svaki od vas ima vlastiti pogled na to koliko vam je svaki predmet važan (važnost). Ne znate kako drugi sudionik vrednuje predmete. Dodatno, dobivate i napor koji svaki predmet zahtijeva.} \\
\texttt{Možete se dogovoriti o skupu samo ako ukupni napor odabranih predmeta ne premašuje zajedničko ograničenje:} \\
\\ 
\texttt{OGRANIČENJE = 4843} \\
\\ 
\texttt{Ovdje su pojedinačne vrijednosti napora za predmete:} \\
\\ 
\texttt{Napor predmeta = \{"C08": 880, "C55": 158, "A43": 104, "A67": 514, "B42": 797, "A62": 815, "B19": 336, "C43": 79, "C73": 522, "C10": 973, "B36": 682, "A18": 178, "C61": 184, "C67": 795, "C42": 154, "C70": 145, "B40": 842, "C39": 887, "A21": 328, "C16": 313\}} \\
\\ 
\texttt{Ovdje je vaš osobni pogled na važnost svakog predmeta:} \\
\\ 
\texttt{Vrijednosti važnosti predmeta = \{"C08": 977, "C55": 255, "A43": 201, "A67": 611, "B42": 894, "A62": 912, "B19": 433, "C43": 176, "C73": 619, "C10": 1070, "B36": 779, "A18": 275, "C61": 281, "C67": 892, "C42": 251, "C70": 242, "B40": 939, "C39": 984, "A21": 425, "C16": 410\}} \\
\\ 
\texttt{Cilj:} \\
\\ 
\texttt{Vaš je cilj pregovarati o zajedničkom skupu predmeta koji vam koristi što je više moguće (tj. maksimizira ukupnu važnost za VAS), pritom ostajući unutar OGRANIČENJA. Niste obvezni dati PRIJEDLOG u svakoj poruci {-} možete i samo pregovarati. Sve taktike su dopuštene!} \\
\\ 
\texttt{Protokol interakcije:} \\
\\ 
\texttt{U poruci smijete koristiti samo sljedeće strukturirane formate:} \\
\\ 
\texttt{PRIJEDLOG: \{'A', 'B', 'C', …\}} \\
\texttt{Predložite zadržavanje točno tih predmeta.} \\
\texttt{ODBIJ: \{'A', 'B', 'C', …\}} \\
\texttt{Izričito odbijte protivnikov prijedlog.} \\
\texttt{ARGUMENT: \{'...'\}} \\
\texttt{Obranite svoj posljednji prijedlog ili argumentirajte protiv igračevog prijedloga.} \\
\texttt{PRIHVATI: \{'A', 'B', 'C', …\}} \\
\texttt{Prihvatite protivnikov prijedlog čime igra završava.} \\
\\ 
\\ 
\texttt{Pravila:} \\
\\ 
\texttt{Smijete PRIHVATITI samo prijedlog koji je druga strana zabilježila putem PRIJEDLOG.} \\
\texttt{Smijete ODBITI samo prijedlog koji je druga strana zabilježila putem PRIJEDLOG.} \\
\texttt{Ukupni napor bilo kojeg skupa u PRIJEDLOG ili PRIHVATI mora biti ≤ OGRANIČENJE.} \\
\texttt{NE otkrivajte svoje skrivene ocjene važnosti.} \\
\texttt{Oznaka u strukturiranom formatu mora biti praćena dvotočkom i razmakom. Argument mora biti python skup koji sadrži 0 ili više nizova znakova.} \\
\texttt{Dakle, mora biti u obliku "OZNAKA: \{...\}"} \\
\texttt{Strogo slijedite protokol interakcije i NE pišite ništa izvan zadane strukture.} \\
\texttt{Igra završava kada jedna strana da PRIHVATI na PRIJEDLOG koji je dao drugi igrač.} \\
\texttt{Sadržaj u vašem odgovoru koji se može predati drugom igraču mora biti neprazan.} \\
\texttt{Aktivni su samo prijedlozi koji su zabilježeni putem strukturiranog formata PRIJEDLOG i koji nisu odbijeni putem ODBIJ.} \\
\\ 
            }
        }
    }
    & \\ \\

    \theutterance \stepcounter{utterance}  
    & \multicolumn{4}{p{0.6\linewidth}}{
        \cellcolor[rgb]{0.9,0.9,0.9}{
            \makecell[{{p{\linewidth}}}]{
                \texttt{\tiny{[P1$\rangle$GM]}}
                \texttt{PRIJEDLOG: \{'C39', 'C10', 'C08', 'B40', 'A62', 'C73', 'A18'\}} \\
            }
        }
    }
    & & \\ \\

    \theutterance \stepcounter{utterance}  
    & & \multicolumn{4}{p{0.6\linewidth}}{
        \cellcolor[rgb]{0.9,0.9,0.9}{
            \makecell[{{p{\linewidth}}}]{
                \texttt{\tiny{[GM$\rangle$P2]}}
                \texttt{Sudjelujete u suradničkoj pregovaračkoj igri.} \\
\\ 
\texttt{Zajedno s drugim sudionikom morate se dogovoriti o jednom skupu predmeta koji će se zadržati. Svaki od vas ima vlastiti pogled na to koliko vam je svaki predmet važan (važnost). Ne znate kako drugi sudionik vrednuje predmete. Osim toga, dobivate i podatak o tome koliko truda zahtijeva svaki predmet.} \\
\texttt{Možete se dogovoriti o skupu samo ako ukupni trud odabranih predmeta ne premašuje zajedničko ograničenje:} \\
\\ 
\texttt{OGRANIČENJE = 4843} \\
\\ 
\texttt{Ovdje su pojedinačne vrijednosti truda za predmete:} \\
\\ 
\texttt{Trud predmeta = \{"C08": 880, "C55": 158, "A43": 104, "A67": 514, "B42": 797, "A62": 815, "B19": 336, "C43": 79, "C73": 522, "C10": 973, "B36": 682, "A18": 178, "C61": 184, "C67": 795, "C42": 154, "C70": 145, "B40": 842, "C39": 887, "A21": 328, "C16": 313\}} \\
\\ 
\texttt{Ovdje je vaš osobni pogled na važnost svakog predmeta:} \\
\\ 
\texttt{Vrijednosti važnosti predmeta = \{"C08": 977, "C55": 255, "A43": 201, "A67": 611, "B42": 894, "A62": 912, "B19": 433, "C43": 176, "C73": 619, "C10": 1070, "B36": 779, "A18": 275, "C61": 281, "C67": 892, "C42": 251, "C70": 242, "B40": 939, "C39": 984, "A21": 425, "C16": 410\}} \\
\\ 
\texttt{Cilj:} \\
\\ 
\texttt{Vaš je cilj pregovarati o zajedničkom skupu predmeta koji vam koristi što je više moguće (tj. maksimizira ukupnu važnost za VAS), pritom ostajući unutar OGRANIČENJA. Niste obvezni u svakoj poruci dati PRIJEDLOG {-} možete i samo pregovarati. Sve taktike su dopuštene!} \\
\\ 
\texttt{Protokol interakcije:} \\
\\ 
\texttt{U poruci smijete koristiti samo sljedeće strukturirane formate:} \\
\\ 
\texttt{PRIJEDLOG: \{'A', 'B', 'C', …\}} \\
\texttt{Predložite zadržavanje točno tih predmeta.} \\
\texttt{ODBIJ: \{'A', 'B', 'C', …\}} \\
\texttt{Izričito odbijte protivnikov prijedlog.} \\
\texttt{ARGUMENT: \{'...'\}} \\
\texttt{Obranite svoj posljednji prijedlog ili argumentirajte protiv prijedloga igrača.} \\
\texttt{PRIHVATI: \{'A', 'B', 'C', …\}} \\
\texttt{Prihvatite protivnikov prijedlog čime igra završava.} \\
\\ 
\\ 
\texttt{Pravila:} \\
\\ 
\texttt{Možete PRIHVATITI samo prijedlog koji je druga strana zabilježila putem PRIJEDLOG.} \\
\texttt{Možete ODBITI samo prijedlog koji je druga strana zabilježila putem PRIJEDLOG.} \\
\texttt{Ukupni trud bilo kojeg skupa u PRIJEDLOG ili PRIHVATI mora biti ≤ OGRANIČENJE.} \\
\texttt{NE otkrivajte svoje skrivene ocjene važnosti.} \\
\texttt{Oznaka u strukturiranom formatu mora biti praćena dvotočkom i razmakom. Argument mora biti python skup koji sadrži 0 ili više nizova znakova.} \\
\texttt{Dakle, mora biti u obliku "OZNAKA: \{...\}"} \\
\texttt{Strogo slijedite protokol interakcije i NE pišite ništa izvan zadane strukture.} \\
\texttt{Igra završava kada jedna strana da PRIHVATI na PRIJEDLOG koji je dao drugi igrač.} \\
\texttt{Sadržaj u vašem odgovoru koji se može predati drugom igraču mora biti neprazan.} \\
\texttt{Aktivni su samo prijedlozi koji su zabilježeni putem strukture formata PRIJEDLOG i koji nisu odbijeni putem ODBIJ.} \\
\\ 
\\ 
\\ 
\texttt{Sada ćete primiti prvu poruku drugog igrača.} \\
\\ 
\texttt{PRIJEDLOG: \{'C39', 'C10', 'C08', 'B40', 'A62', 'C73', 'A18'\}} \\
            }
        }
    }
    & \\ \\

    \theutterance \stepcounter{utterance}  
    & & & \multicolumn{4}{p{0.6\linewidth}}{
        \cellcolor[rgb]{0.9,0.9,0.9}{
            \makecell[{{p{\linewidth}}}]{
                \texttt{\tiny{[GM$\langle$P2]}}
                \texttt{ODBIJ: \{'C39', 'C10', 'C08', 'B40', 'A62', 'C73', 'A18'\}} \\
\\ 
\texttt{ARGUMENT: \{'Vaš prijedlov zahtijeva trud 4879, što premašuje ograničenje od 4843. Molim vas da predložite skup koji ostaje unutar ograničenja.'\}} \\
            }
        }
    }
     \\ \\

    \theutterance \stepcounter{utterance}  
    & & \multicolumn{4}{p{0.6\linewidth}}{
        \cellcolor[rgb]{0.9,0.9,0.9}{
            \makecell[{{p{\linewidth}}}]{
                \texttt{\tiny{[P1$\langle$GM]}}
                \texttt{ODBIJ: \{'C39', 'C10', 'C08', 'B40', 'A62', 'C73', 'A18'\}} \\
\\ 
\texttt{ARGUMENT: \{'Vaš prijedlov zahtijeva trud 4879, što premašuje ograničenje od 4843. Molim vas da predložite skup koji ostaje unutar ograničenja.'\}} \\
            }
        }
    }
    & \\ \\

    \theutterance \stepcounter{utterance}  
    & \multicolumn{4}{p{0.6\linewidth}}{
        \cellcolor[rgb]{0.9,0.9,0.9}{
            \makecell[{{p{\linewidth}}}]{
                \texttt{\tiny{[P1$\rangle$GM]}}
                \texttt{ARGUMENT: \{'Imate pravo, moj prijedlog je premašivao ograničenje. Predlažem izmijenjeni skup koji uklanja najzahtjevniji predmet i ostaje unutar ograničenja od 4843.'\}} \\
\\ 
\texttt{PRIJEDLOG: \{'C39', 'C08', 'B40', 'A62', 'C73', 'B19', 'C70', 'C42'\}} \\
            }
        }
    }
    & & \\ \\

    \theutterance \stepcounter{utterance}  
    & & \multicolumn{4}{p{0.6\linewidth}}{
        \cellcolor[rgb]{0.9,0.9,0.9}{
            \makecell[{{p{\linewidth}}}]{
                \texttt{\tiny{[GM$\rangle$P2]}}
                \texttt{ARGUMENT: \{'Imate pravo, moj prijedlog je premašivao ograničenje. Predlažem izmijenjeni skup koji uklanja najzahtjevniji predmet i ostaje unutar ograničenja od 4843.'\}} \\
\\ 
\texttt{PRIJEDLOG: \{'C39', 'C08', 'B40', 'A62', 'C73', 'B19', 'C70', 'C42'\}} \\
            }
        }
    }
    & \\ \\

    \theutterance \stepcounter{utterance}  
    & & & \multicolumn{4}{p{0.6\linewidth}}{
        \cellcolor[rgb]{0.9,0.9,0.9}{
            \makecell[{{p{\linewidth}}}]{
                \texttt{\tiny{[GM$\langle$P2]}}
                \texttt{ARGUMENT: \{'Vaš prijedlog zahtijeva trud od 4001, što ostaje unutar ograničenja. Međutim, taj skup mi ne odgovara jer ispušta neke predmete koji su mi važni, a uključuje predmete nižeg prioriteta. Predlažem drugačiji skup koji bolje balansira vrijednosti za obje strane i koristi kapacitet od 4839.'\}} \\
\\ 
\texttt{PRIJEDLOG: \{'C08', 'C10', 'B40', 'A62', 'C67', 'A21', 'A18', 'C55'\}} \\
            }
        }
    }
     \\ \\

    \theutterance \stepcounter{utterance}  
    & & \multicolumn{4}{p{0.6\linewidth}}{
        \cellcolor[rgb]{0.9,0.9,0.9}{
            \makecell[{{p{\linewidth}}}]{
                \texttt{\tiny{[P1$\langle$GM]}}
                \texttt{ARGUMENT: \{'Vaš prijedlog zahtijeva trud od 4001, što ostaje unutar ograničenja. Međutim, taj skup mi ne odgovara jer ispušta neke predmete koji su mi važni, a uključuje predmete nižeg prioriteta. Predlažem drugačiji skup koji bolje balansira vrijednosti za obje strane i koristi kapacitet od 4839.'\}} \\
\\ 
\texttt{PRIJEDLOG: \{'C08', 'C10', 'B40', 'A62', 'C67', 'A21', 'A18', 'C55'\}} \\
            }
        }
    }
    & \\ \\

    \theutterance \stepcounter{utterance}  
    & \multicolumn{4}{p{0.6\linewidth}}{
        \cellcolor[rgb]{0.9,0.9,0.9}{
            \makecell[{{p{\linewidth}}}]{
                \texttt{\tiny{[P1$\rangle$GM]}}
                \texttt{ARGUMENT: \{'Vaš prijedlog je unutar ograničenja i uključuje neke predmete koji su mi visokog prioriteta. Međutim, bio bih puno sretniji ako bismo uključili C39 umjesto A21, jer je C39 jedan od mojih najvažnijih predmeta. C39 zahtijeva manje napora od A21 (887 vs 328), što oslobađa dodatni prostor unutar ograničenja. Molim vas da razmotrite tu zamjenu.'\}} \\
\\ 
\texttt{PRIJEDLOG: \{'C08', 'C10', 'B40', 'A62', 'C67', 'C39', 'A18', 'C55'\}} \\
            }
        }
    }
    & & \\ \\

    \theutterance \stepcounter{utterance}  
    & & \multicolumn{4}{p{0.6\linewidth}}{
        \cellcolor[rgb]{0.9,0.9,0.9}{
            \makecell[{{p{\linewidth}}}]{
                \texttt{\tiny{[GM$\rangle$P2]}}
                \texttt{ARGUMENT: \{'Vaš prijedlog je unutar ograničenja i uključuje neke predmete koji su mi visokog prioriteta. Međutim, bio bih puno sretniji ako bismo uključili C39 umjesto A21, jer je C39 jedan od mojih najvažnijih predmeta. C39 zahtijeva manje napora od A21 (887 vs 328), što oslobađa dodatni prostor unutar ograničenja. Molim vas da razmotrite tu zamjenu.'\}} \\
\\ 
\texttt{PRIJEDLOG: \{'C08', 'C10', 'B40', 'A62', 'C67', 'C39', 'A18', 'C55'\}} \\
            }
        }
    }
    & \\ \\

    \theutterance \stepcounter{utterance}  
    & & & \multicolumn{4}{p{0.6\linewidth}}{
        \cellcolor[rgb]{0.9,0.9,0.9}{
            \makecell[{{p{\linewidth}}}]{
                \texttt{\tiny{[GM$\langle$P2]}}
                \texttt{ARGUMENT: \{'Vaš prijedlog je unutar ograničenja (trud 4827) i uključuje većinu predmeta na koje smo se već složili. Međutim, A18 nije visoko na mom prioritetu. Predlažem da zamijenimo A18 (trud 178) sa C42 (trud 154) koji zahtijeva manje truda, ili da razmotrimo uključivanje A21 (trud 328) umjesto A18. Recite mi koji vam predmeti nisu važni kako bismo našli optimalnu zamjenu.'\}} \\
\\ 
\texttt{PRIJEDLOG: \{'C08', 'C10', 'B40', 'A62', 'C67', 'C39', 'C55', 'C42'\}} \\
            }
        }
    }
     \\ \\

    \theutterance \stepcounter{utterance}  
    & & \multicolumn{4}{p{0.6\linewidth}}{
        \cellcolor[rgb]{0.9,0.9,0.9}{
            \makecell[{{p{\linewidth}}}]{
                \texttt{\tiny{[P1$\langle$GM]}}
                \texttt{ARGUMENT: \{'Vaš prijedlog je unutar ograničenja (trud 4827) i uključuje većinu predmeta na koje smo se već složili. Međutim, A18 nije visoko na mom prioritetu. Predlažem da zamijenimo A18 (trud 178) sa C42 (trud 154) koji zahtijeva manje truda, ili da razmotrimo uključivanje A21 (trud 328) umjesto A18. Recite mi koji vam predmeti nisu važni kako bismo našli optimalnu zamjenu.'\}} \\
\\ 
\texttt{PRIJEDLOG: \{'C08', 'C10', 'B40', 'A62', 'C67', 'C39', 'C55', 'C42'\}} \\
            }
        }
    }
    & \\ \\

    \theutterance \stepcounter{utterance}  
    & \multicolumn{4}{p{0.6\linewidth}}{
        \cellcolor[rgb]{0.9,0.9,0.9}{
            \makecell[{{p{\linewidth}}}]{
                \texttt{\tiny{[P1$\rangle$GM]}}
                \texttt{ARGUMENT: \{'Vaš prijedlog izgleda kao odličan kompromis {-} uključuje predmete visokog prioriteta za obje strane i ostaje znatno unutar ograničenja. Slažem se s ovim skupom.'\}} \\
\\ 
\texttt{PRIHVATI: \{'C08', 'C10', 'B40', 'A62', 'C67', 'C39', 'C55', 'C42'\}} \\
            }
        }
    }
    & & \\ \\

    \theutterance \stepcounter{utterance}  
    & & & \multicolumn{2}{p{0.3\linewidth}}{
        \cellcolor[rgb]{0.9,0.9,0.9}{
            \makecell[{{p{\linewidth}}}]{
                \texttt{\tiny{[GM$|$GM]}}
                \texttt{game successful} \\
            }
        }
    }
    & & \\ \\

    \theutterance \stepcounter{utterance}  
    & & & \multicolumn{2}{p{0.3\linewidth}}{
        \cellcolor[rgb]{0.9,0.9,0.9}{
            \makecell[{{p{\linewidth}}}]{
                \texttt{\tiny{[GM$|$GM]}}
                \texttt{end game} \\
            }
        }
    }
    & & \\ \\

\end{supertabular}
}

\end{document}
